\section{Detector error model}
\label{sec:detector-error-model}

We use the detector error model (DEM)~\cite{gidney_stim_2021} as a compact description of the noisy syndrome extraction circuit.
For our purposes, the DEM serves two roles.
First, it provides a convenient first-order language for identifying harmful error mechanisms that survive postselection and induce logical injection errors.
Second, it yields a simple explanatory proxy for the discard rate, which is one of the main quantities entering the spacetime-overhead analysis.
This appendix also explains why the discard rate remains similar across different optimized initial configurations.

\subsection{Definitions}

A \emph{detector} is a binary parity check built from a specified subset of measurement-record bits.
Detectors are defined so that, in the absence of noise, every detector outcome is deterministically $0$.
Noise can toggle some of these detector outcomes to $1$.

A DEM records how elementary stochastic mechanisms toggle detector outcomes and tracked logical observables, without explicitly tracking the full quantum state.
For the analysis below, we represent the DEM as a collection of Bernoulli mechanisms indexed by $e$.
Each mechanism is specified by a triple
\begin{equation}
(p_e, D_e, L_e),
\label{eq:dem_mechanism_triple}
\end{equation}
where $p_e\in[0,1]$ is the occurrence probability of mechanism $e$, $D_e\subseteq\{1,\dots,N_D\}$ is the set of detectors toggled by $e$, and $L_e\subseteq\{1,\dots,N_L\}$ is the set of tracked logical observables toggled by $e$.
Here $N_L$ is the number of logical observables included in the DEM analysis, and $N_D$ is the number of detectors.

\subsection{Error mechanism counting}

For MSI, the tracked logical observables are the injected logical observables indexed by $\cI$, so we take $L_e\subseteq \cI$.
The most dangerous first-order mechanisms are those with
\begin{equation}
L_e \neq \emptyset,
\qquad
D_e = \emptyset,
\end{equation}
because they flip one or more injected logical observables without triggering any detector.
Such mechanisms evade postselection and also leave no detector information for decoding.

Because detector and logical outcomes combine by parity, multiple mechanisms can cancel each other.
However, the probability that two distinct mechanisms $e_1$ and $e_2$ occur in the same shot is second order, namely $p_{e_1}p_{e_2}$.
Accordingly, a standard first-order approximation is to assume that at most one DEM mechanism occurs during a shot.
Previous MSI analyses use this viewpoint to count undetected logical-fault mechanisms and find good agreement with circuit-level Monte Carlo simulation~\cite{li_magic_2015,lao_magic_2022,gidney_cleaner_2023}.
In our setting, this yields
\begin{equation}
p_\cJ \approx \sum_{e:\,L_e = \cJ} p_e,
\end{equation}
where $p_\cJ$ is the probability that the incorrect logicals are \emph{exactly} the subset $\cJ$.
In \Cref{sec:evaluation}, we show that this first-order approximation closely matches circuit-level Monte Carlo simulation.

\subsection{Postselection and discard rate}

Let $P\subseteq\{1,\dots,N_D\}$ be the set of postselection detectors.
In our simulations, $P$ consists of (i) fixed-stabilizer detectors from the first SE round and (ii) round-parity detectors comparing the first two SE rounds.
A shot is \emph{accepted} if and only if all detectors in $P$ are $0$; otherwise the shot is discarded.

Define the set of \emph{postselection-relevant} mechanisms by
\begin{equation}
R := \{e:\; D_e\cap P\neq \emptyset\}.
\label{eq:relevant_mechanisms_set_simplified}
\end{equation}
If no mechanism in $R$ occurs, then every postselection detector remains $0$.
Therefore, discard implies that at least one mechanism in $R$ occurred, which gives the upper bound
\begin{equation}
r_{\rm disc}
\le
\Pr\!\left(\exists e\in R:\; e\text{ occurs}\right).
\end{equation}
The inequality above is generally strict because multiple postselection-relevant mechanisms can occur simultaneously and cancel each other by parity.

Under the Bernoulli-mechanism interpretation above, the probability that no mechanism in $R$ occurs is
\begin{equation}
a := \prod_{e\in R}(1-p_e).
\end{equation}
We call $a$ the \emph{acceptance proxy}.
Equivalently,
\begin{equation}
\Pr\!\left(\exists e\in R:\; e\text{ occurs}\right)
= 1-\prod_{e\in R}(1-p_e)
= 1-a.
\end{equation}
This motivates the \emph{union proxy} for the discard rate,
\begin{equation}
r_{\rm disc}^{\rm UP} := 1-a = 1-\prod_{e\in R}(1-p_e).
\label{eq:rdisc_up_def}
\end{equation}
As shown in \Cref{sec:evaluation}, the union proxy closely tracks the discard rate measured by circuit-level Monte Carlo simulation, $r_{\rm disc}^{\rm MC}$.

\subsection{Acceptance factor decomposition}

To explain why the discard rate varies little across optimized initial configurations, it is useful to separate mechanisms caught immediately by fixed-stabilizer detectors from those caught only by round-parity detectors.
Let $P_{\rm fix}\subseteq P$ denote the set of fixed-stabilizer detectors in the first SE round, and let $|F|:=|P_{\rm fix}|$ be the number of fixed stabilizers.
Then the acceptance proxy factorizes as
\begin{align}
a
&= \prod_{e:\,D_e\cap P\neq \emptyset}(1-p_e) \\
&=
\underbrace{\prod_{e:\,D_e\cap P_{\rm fix}\neq \emptyset}(1-p_e)}_{a_{\rm fix}}
\cdot
\underbrace{\prod_{\substack{e:\,D_e\cap P\neq \emptyset\\ D_e\cap P_{\rm fix}=\emptyset}}(1-p_e)}_{a_{{\rm par}\mid{\rm fix}}}.
\label{eq:acceptance_decomp}
\end{align}
Here $a_{\rm fix}$ captures the mechanisms that would be flagged directly by fixed-stabilizer detectors, while $a_{{\rm par}\mid{\rm fix}}$ captures the remaining postselection-relevant mechanisms that do \emph{not} trigger any fixed-stabilizer detector and are therefore caught only by round-parity detectors.
The two underlying mechanism sets are disjoint, so the product factorization is exact.
Substituting \Cref{eq:acceptance_decomp} into \Cref{eq:rdisc_up_def} gives
\begin{equation}
r_{\rm disc}^{\rm UP}(P)
=
1-a_{\rm fix}a_{{\rm par}\mid{\rm fix}}.
\label{eq:rdisc_up_factored}
\end{equation}

\subsection{Explanation of similar discard rate across different initial configurations}
\label{sec:explanation-similar-discard-rate}

\begin{table}[htb]
    \centering
    \begin{tabular}{ccccccc}
    \toprule
    $p$ & $|\mathcal{I}|$ & $|F|$    & $(a_{\rm fix},\, a_{{\rm par}\mid{\rm fix}})$ & $r_{\rm disc}^{\rm UP}(p)$ & $r_{\rm disc}^{\rm MC}(p)$ \\
    \midrule
$10^{-4}$ & 1 & 26 & (0.9563, 0.8018) & 0.2333 & 0.2331  \\
$10^{-4}$ & 2 & 32 & (0.9487, 0.8062) & 0.2351 & 0.2352  \\
$10^{-4}$ & 3 & 31 & (0.9495, 0.8056) & 0.2351 & 0.2346  \\
$10^{-4}$ & 4 & 32 & (0.9484, 0.8061) & 0.2355 & 0.2352  \\
$10^{-4}$ & 5 & 28 & (0.9517, 0.8040) & 0.2348 & 0.2346  \\
$10^{-4}$ & 6 & 18 & (0.9650, 0.7965) & 0.2314 & 0.2315  \\
\midrule
$5 \times 10^{-4}$ & 1 & 26 & (0.7996, 0.3313) & 0.7351 & 0.7344  \\
$5 \times 10^{-4}$ & 2 & 32 & (0.7683, 0.3407) & 0.7383 & 0.7380  \\
$5 \times 10^{-4}$ & 3 & 31 & (0.7716, 0.3393) & 0.7382 & 0.7383  \\
$5 \times 10^{-4}$ & 4 & 32 & (0.7672, 0.3404) & 0.7388 & 0.7387  \\
$5 \times 10^{-4}$ & 5 & 28 & (0.7808, 0.3360) & 0.7376 & 0.7382  \\
$5 \times 10^{-4}$ & 6 & 18 & (0.8367, 0.3205) & 0.7318 & 0.7316  \\
\midrule
$10^{-3}$ & 1 & 26 & (0.6394, 0.1097) & 0.9298 & 0.9300  \\
$10^{-3}$ & 2 & 32 & (0.5903, 0.1160) & 0.9315 & 0.9312  \\
$10^{-3}$ & 3 & 31 & (0.5954, 0.1151) & 0.9315 & 0.9311  \\
$10^{-3}$ & 4 & 32 & (0.5886, 0.1159) & 0.9318 & 0.9318  \\
$10^{-3}$ & 5 & 28 & (0.6096, 0.1129) & 0.9312 & 0.9313  \\
$10^{-3}$ & 6 & 18 & (0.7001, 0.1027) & 0.9281 & 0.9278  \\
    \bottomrule
    \end{tabular}
    \caption{For each physical error rate $p$ and injection configuration $|\mathcal{I}|$, we report the Monte Carlo discard rate $r_{\rm disc}^{\rm MC}$, the union proxy $r_{\rm disc}^{\rm UP}$, and the two acceptance factors $a_{\rm fix}$ and $a_{{\rm par}\mid{\rm fix}}$ defined in \Cref{eq:acceptance_decomp}. The factorization \(r_{\rm disc}^{\rm UP}=1-a_{\rm fix}a_{{\rm par}\mid{\rm fix}}\) is given by \Cref{eq:rdisc_up_factored}.}
    \label{tab:discard_components}
\end{table}


\Cref{tab:discard_components} shows that the union proxy tracks the Monte Carlo discard rate very closely for all optimized BB code configurations we tested.
It also explains why the discard rate changes only weakly across these configurations, even though the number of fixed stabilizers $|F|$ varies substantially.
From
\begin{equation}
    r_{\rm disc}^{\rm UP} = 1 - a_{\rm fix} a_{{\rm par}\mid{\rm fix}},
\end{equation}
we obtain
\begin{equation}
    \frac{\partial r_{\rm disc}^{\rm UP}}{\partial a_{\rm fix}} = -a_{{\rm par}\mid{\rm fix}}.
\end{equation}
Thus, for a given change $\Delta a_{\rm fix}$,
\begin{equation}
|\Delta r_{\rm disc}^{\rm UP}| \approx a_{{\rm par}\mid{\rm fix}}\,|\Delta a_{\rm fix}|,
\end{equation}
so the sensitivity of the discard rate to changes in $a_{\rm fix}$ is itself controlled by $a_{{\rm par}\mid{\rm fix}}$.

A representative comparison is provided by the optimized configurations with $|\cI|=5$ and $|\cI|=6$.
At high physical error rate, for example $p=10^{-3}$, the change in $a_{\rm fix}$ is relatively large:
\begin{equation}
\Delta a_{\rm fix} \approx 0.7001 - 0.6096 \approx 0.091.
\end{equation}
However, in this regime $a_{{\rm par}\mid{\rm fix}}$ is small, about $0.10$, so the resulting change in discard rate is close to $10^{-2}$.
By contrast, at low physical error rate, for example $p=10^{-4}$, the coefficient $a_{{\rm par}\mid{\rm fix}}$ is much larger, about $0.80$, but the change in $a_{\rm fix}$ becomes much smaller:
\begin{equation}
\Delta a_{\rm fix} \approx 0.9650 - 0.9517 \approx 0.013.
\end{equation}
Again, the resulting change in discard rate remains only of order $10^{-2}$.

The same qualitative mechanism appears across the optimized configurations in \Cref{tab:discard_components}.
At larger $p$, round-parity-only mechanisms are so prevalent that $a_{{\rm par}\mid{\rm fix}}$ is small, which suppresses the impact of variations in $a_{\rm fix}$.
At smaller $p$, $a_{{\rm par}\mid{\rm fix}}$ is larger, but the optimized configurations have much closer values of $a_{\rm fix}$.
Therefore, even substantial changes in the number of fixed stabilizers do not translate into equally large changes in the total discard rate.
This explains why the discard rate remains similar across the optimized initial configurations considered in \Cref{sec:evaluation}.
